%%
%% fall-lecture16.tex
%% 
%% Made by Alex Nelson <pqnelson@gmail.com>
%% Login   <alex@lisp>
%% 
%% Started on  2026-11-03T08:43:05-0800
%% Last update 2026-11-03T08:43:05-0800
%% 

\lecture{}

\subsection{Recipricocity Relations}

\begin{node}
We want to have 2 equations we can compare side-by-side.
\end{node}

\begin{example}
Enthalpy has a state equation of the form:
\begin{equation}
\D h = T\,\D S + V\,\D p.
\end{equation}
This implies
\begin{equation}
T\,\D S = \D h - V\,\D p,
\end{equation}
giving us
\begin{equation}\label{eq:fall:lec16:ex1:eq1}
\D S = \frac{\D h}{T}-\frac{V}{T}\,\D p.
\end{equation}
This is just the first step.

Now, we also can write enthalpy as a function of temperature and
pressure $h=h(T,p)$ which gives us the difference equation
\begin{equation}
\D h = \left(\frac{\partial h}{\partial p}\right)_{T}\D p
+ \left(\frac{\partial h}{\partial T}\right)_{p}\D T
\end{equation}
Then plugging this back into Eq~\eqref{eq:fall:lec16:ex1:eq1},
\begin{subequations}\label{eq:fall:lec16:ex1:eq3}
  \begin{align}
\D S &= \frac{1}{T}\left[\left(\frac{\partial h}{\partial p}\right)_{T}\D p
+ \left(\frac{\partial h}{\partial T}\right)_{p}\D T\right]-\frac{V}{T}\D p\\
&=\left[\frac{1}{T}\left(\frac{\partial h}{\partial p}\right)_{T}
 -\frac{V}{T} \right]\D p + \left(\frac{\partial h}{\partial T}\right)_{p}\D T
  \end{align}
\end{subequations}
This is one equation.

Now, we can write entropy as a function of pressure and temperature $S=S(p,T)$
which gives us the difference equation
\begin{equation}\label{eq:fall:lec16:ex1:eq4}
\D S = \left(\frac{\partial S}{\partial p}\right)_{T}\D p
+ \left(\frac{\partial S}{\partial T}\right)_{p}\D T.
\end{equation}
This is our second equation. Now setting Eq~\eqref{eq:fall:lec16:ex1:eq3}
equal to Eq~\eqref{eq:fall:lec16:ex1:eq4} we obtain (matching
coefficients of the differentials) two equations:
\begin{subequations}
\begin{equation}
\left(\frac{\partial S}{\partial p}\right)_{T} = \frac{1}{T}\left(\frac{\partial h}{\partial p}\right)_{T}
\end{equation}
and
\begin{equation}
\left(\frac{\partial S}{\partial T}\right)_{p}
=\frac{1}{T}\left(\frac{\partial h}{\partial T}\right)_{p}-\frac{V}{T}.
\end{equation}
\end{subequations}
We're not done yet! We can rearrange this second equation and the
Maxwell relation for Gibbs free energy Eq~\eqref{eq:maxwell-relation-for-gibbs-free-energy} to obtain a
\define{Reciprocity Relation}
\begin{equation}
\boxed{\left(\frac{\partial h}{\partial T}\right)_{p}=V-T\left(\frac{\partial S}{\partial T}\right)_{p}.}
\end{equation}
\end{example}

\begin{node}
There are only 2 reciprocity relations. The other uses the canonical
equation of state
\begin{equation}
\D e = T\,\D S - p\,\D V,
\end{equation}
and writes entropy as a function of volume and temperature $S=S(V,T)$
to obtain
\begin{subequations}
\begin{equation}
  \left(\frac{\partial S}{\partial T}\right)_{V}=
  \frac{1}{T}\left(\frac{\partial e}{\partial T}\right)_{V}
\end{equation}
and
\begin{equation}
\left(\frac{\partial S}{\partial V}\right)_{T} =
\frac{1}{T}\left[\left(\frac{\partial e}{\partial V}\right)_{T} + p\right]
\end{equation}
\end{subequations}
Then after a lot of profanity and invoking the Helmholtz free energy
Maxwell relation Eq~\eqref{eq:maxwell-relation-for-helmholtz-free-energy}, the
second equation gives us the \define{Second Reciprocity Relation}:
\begin{equation}
\boxed{\left(\frac{\partial e}{\partial V}\right)_{T} = T\left(\frac{\partial p}{\partial T}\right)_{V}-p.}
\end{equation}
(This is just the other reciprocity relation, ``second'' is used to
distinguish it from the earlier relation we just so-happend to derive first.)
\end{node}

\begin{remark}
There were more than a few intermediate steps involved obtaining these
reciprocity relations, and they are all valid.
\end{remark}

\begin{remark}
I am having trouble finding anyone using the term ``reciprocity relation''
for either of these two equations. This just means I am not competent
enough to find references.
\end{remark}

\subsection{Defined quantities in Thermodynamics}

\begin{definition}
The \define{Specific heat at constant pressure} is
\begin{equation*}
c_{p} := \left(\frac{\partial h}{\partial T}\right)_{p} = T\left(\frac{\partial S}{\partial T}\right)_{p}
\end{equation*}
where we use the equation of state to get the equality with $\partial S/\partial T$.
\end{definition}

\begin{definition}
The \define{Specific Heat at Constant Volume} is
\begin{equation*}
c_{V} := \left(\frac{\partial h}{\partial T}\right)_{V} = T\left(\frac{\partial S}{\partial T}\right)_{V}
\end{equation*}
where we use an equation of state to get the equality with $\partial S/\partial T$.
\end{definition}

\begin{remark}
We can think of the specific heats as functions $c_{p}=c_{p}(p,T)$ and $c_{V}=c_{V}(V,T)$.
\end{remark}

\begin{definition}
The \define{Ratio of Specific Heats} is $\gamma := c_{p}/c_{V}$.
\end{definition}

\begin{definition}
The \define{Speed of Sound} (as a thermodynamic quantity) is
\begin{equation*}
c := \sqrt{\left(\frac{\partial p}{\partial \rho}\right)_{S}}
\end{equation*}
where $\rho=1/V$. We should think of this as describing compressive
forces (how pressure changes with respect to specific volume).
\end{definition}

\begin{definition}
The \define{Coefficient of Thermal Expansion} is
\begin{equation*}
\alpha := \frac{1}{V}\left(\frac{\partial V}{\partial T}\right)_{p}.
\end{equation*}
Typically $\alpha>0$, but famously $\alpha<0$ for water between
$0^{\circ}~\mathrm{C}$ to $4^{\circ}~\mathrm{C}$.
\end{definition}

\begin{definition}
We define the \define{Isothermal Compressibility}
\begin{equation}
\kappa_{T} := \frac{1}{V}\left(\frac{\partial V}{\partial p}\right)_{T}
\end{equation}
and \define{Isentropic Compressibility}
\begin{equation}
\kappa_{S} := \frac{-1}{V}\left(\frac{\partial V}{\partial p}\right)_{S}.
\end{equation}
The convention in Thermodynamics is to reserve kappa $\kappa$ for
``compressibility'' quantities.
\end{definition}

\begin{remark}
Note that
\begin{subequations}
  \begin{align}
\left(\frac{\partial V}{\partial p}\right)_{S}
&= \frac{1}{\left(\frac{\partial p}{\partial V}\right)_{S}}\\
&= \frac{1}{\left(\frac{\partial p}{\partial \rho}\right)_{S}\left(\frac{\partial \rho}{\partial V}\right)_{S}}\\
&= \frac{1}{c^{2}\frac{-1}{V^{2}}},
  \end{align}
\end{subequations}
so
\begin{equation}
\kappa_{S} = \frac{V}{c^{2}}
\end{equation}
where $c$ is the speed of sound.
\end{remark}

\begin{proposition}
Isothermal compressibility is directly proportional to isentropic
compressibility, with the coefficient of proportionality being
precisely the ratio of specific heats:
\begin{equation}
\kappa_{T} = \gamma\kappa_{S}.
\end{equation}
\end{proposition}

(See, e.g., Landau and Lifshitz~\cite[\S16]{landau1969statmech}
specifically their derivation for Eq~(16.14).)

\begin{proof}
Observe we want to prove
\begin{equation}
\left(\frac{\partial p}{\partial V}\right)_{S} = \frac{c_{p}}{c_{V}}\left(\frac{\partial p}{\partial V}\right)_{T}.
\end{equation}
We start by writing entropy as a function of pressure and volume $S=S(p,V)$.
Then we write the difference equation
\begin{equation}
\D S = \left(\frac{\partial S}{\partial p}\right)_{V}\D p
+ \left(\frac{\partial S}{\partial V}\right)_{p}\D V.
\end{equation}
Now, looking at
\begin{subequations}
  \begin{align}
\left(\frac{\partial S}{\partial p}\right)_{V} &= \left(\frac{\partial S}{\partial T}\right)_{V}\left(\frac{\partial T}{\partial p}\right)_{V}\\
&=\frac{c_{V}}{T}\left(\frac{\partial T}{\partial p}\right)_{V}
  \end{align}
\end{subequations}
and
\begin{subequations}
  \begin{align}
\left(\frac{\partial S}{\partial V}\right)_{p}
&= \left(\frac{\partial S}{\partial T}\right)_{p}\left(\frac{\partial T}{\partial V}\right)_{p}\\
&= \frac{c_{P}}{T}\left(\frac{\partial T}{\partial V}\right)_{p},
  \end{align}
\end{subequations}
we can plug these into the difference equation to obtain
\begin{subequations}\label{eq:fall:lec16:difference-spec-heat:step3}
  \begin{align}
\D S &=
\frac{c_{V}}{T}\left(\frac{\partial T}{\partial p}\right)_{V}\D p
+ \frac{c_{P}}{T}\left(\frac{\partial T}{\partial V}\right)_{p}\D V\\
&=\frac{c_{V}}{T}\frac{1}{\left(\frac{\partial p}{\partial T}\right)_{V}}\D p
+ \frac{c_{P}}{T}\frac{1}{\left(\frac{\partial V}{\partial T}\right)_{p}}\D V.
  \end{align}
\end{subequations}
Now, we just need to calculate $\partial p/\partial T$ and $\partial V/\partial T$.

We can write the specific volume as a function of pressure and
temperature $V=V(p,T)$. Then we have the difference equation
\begin{equation}
\D V = \left(\frac{\partial V}{\partial p}\right)_{T}\D p +  \left(\frac{\partial V}{\partial T}\right)_{p}\D T.
\end{equation}
If we divide both sides by $\D T$ and hold $V$ constant, we see
\begin{equation}
 \underbrace{\left(\frac{\partial V}{\partial T}\right)_{V}}_{=0}
= \left(\frac{\partial V}{\partial p}\right)_{T}\left(\frac{\partial p}{\partial T}\right)_{V} +  \left(\frac{\partial V}{\partial T}\right)_{p}.
\end{equation}
Then rearranging terms gives us (since the left-hand side vanishes):
\begin{equation}\label{eq:fall:lec16:difference-spec-heat:step5}
\left(\frac{\partial p}{\partial T}\right)_{V}
 = \frac{-\left(\frac{\partial V}{\partial T}\right)_{p}}{\left(\frac{\partial V}{\partial p}\right)_{T}}=\frac{\alpha}{\kappa_{T}}.
\end{equation}
Similarly, we find
\begin{equation}\label{eq:fall:lec16:difference-spec-heat:step6}
\left(\frac{\partial V}{\partial T}\right)_{p}
= V\alpha.
\end{equation}
Great.

Now, we plug Eq~\eqref{eq:fall:lec16:difference-spec-heat:step5} and Eq~\eqref{eq:fall:lec16:difference-spec-heat:step6}
back into Eq~\eqref{eq:fall:lec16:difference-spec-heat:step3}
to get
\begin{equation}
\D S = \frac{c_{V}}{T}\frac{\kappa_{T}}{\alpha}\D p + \frac{c_{P}}{T}\frac{1}{V\alpha}\D V.
\end{equation}
Performing a similar trick (dividing by $\D V$ while holding $S$
constant) gives us
\begin{equation}
0 = \frac{c_{V}}{T}\frac{\kappa_{T}}{\alpha}\left(\frac{\partial p}{\partial V}\right)_{S} + \frac{c_{P}}{T}\frac{1}{V\alpha}.
\end{equation}
Hence the claim after a bit of algebra.
\end{proof}

\begin{remark}
Landau and Lifshitz take a different approach, which is worth
mentioning. They use the Jacobian matrix determinant
\begin{equation}
\frac{\partial(u,v)}{\partial(x,y)}=\det\begin{pmatrix}\partial
u/\partial x & \partial u/\partial y\\
\partial v/\partial x & \partial v/\partial y
\end{pmatrix}.
\end{equation}
This has the identities
\begin{subequations}
\begin{equation}
\frac{\partial(u,y)}{\partial(x,y)}=\left(\frac{\partial u}{\partial x}\right)_{y},
\end{equation}
and
\begin{equation}
\frac{\partial(u,v)}{\partial(x,y)}=\frac{[\partial(u,v)/\partial(a,b)]}{[\partial(x,y)/\partial(a,b)]}.
\end{equation}
\end{subequations}
Then
\begin{subequations}
  \begin{align}
\left(\frac{\partial T}{\partial V}\right)_{S}
&= \frac{\partial(T,S)}{\partial(V,S)}\\
&= \frac{[\partial(T,S)/\partial(V,T)]}{[\partial(V,S)/\partial(V,T)]}\\
&= -\frac{\left(\frac{\partial S}{\partial V}\right)_{T}}{\left(\frac{\partial S}{\partial T}\right)_{V}}\\
&= \frac{-T}{c_{V}}{\left(\frac{\partial S}{\partial V}\right)_{T}},
  \end{align}
\end{subequations}
and by using the Maxwell relation for Helmholtz free energy rewrites the
right-hand side to give us:
\begin{equation}
\left(\frac{\partial T}{\partial V}\right)_{S} = \frac{T}{c_{V}}\left(\frac{\partial p}{\partial T}\right)_{V}.
\end{equation}
``Similarly, we find'' (using the Maxwell relation for Gibbs free energy),
\begin{equation}
\left(\frac{\partial T}{\partial p}\right)_{S} = \frac{T}{c_{p}}\left(\frac{\partial V}{\partial T}\right)_{p}.
\end{equation}
Then using Jacobian matrix identities again,
\begin{subequations}
  \begin{align}
\left(\frac{\partial V}{\partial p}\right)_{S}
&= \frac{\partial(V,S)}{\partial(p,S)}\\
&= \frac{[\partial(V,S)/\partial(V,T)]}{[\partial(p,S)/\partial(p,T)]}\frac{\partial(V,T)}{\partial(p,T)}\\
&=\frac{\left(\frac{\partial S}{\partial T}\right)_{V}}{\left(\frac{\partial S}{\partial T}\right)_{p}}\left(\frac{\partial V}{\partial p}\right)_{T}\\
&=\frac{c_{V}}{c_{p}}\left(\frac{\partial V}{\partial p}\right)_{T}
  \end{align}
\end{subequations}
then we have the desired result.
\end{remark}

\begin{proposition}
The difference in specific heats satisfies
\begin{equation}
c_{p}-c_{V} = \frac{VT\alpha^{2}}{\kappa_{T}}.
\end{equation}
\end{proposition}

\begin{proof}
Landau and Lifshitz observes
\begin{subequations}
  \begin{align}
c_{V} &= T\left(\frac{\partial S}{\partial T}\right)_{V}\\
&= T \frac{\partial(S,V)}{\partial(T,V)}\\
&= T \frac{[\partial(S,V)/\partial(T,p)]}{[\partial(T,V)/\partial(T,p)]}\\
&= T \frac{\left(\frac{\partial S}{\partial T}\right)_{p}\left(\frac{\partial V}{\partial p}\right)_{T}-\left(\frac{\partial S}{\partial p}\right)_{T}\left(\frac{\partial V}{\partial T}\right)_{p}}{\left(\frac{\partial V}{\partial p}\right)_{T}}\\
&= c_{p} - T \frac{\left(\frac{\partial S}{\partial p}\right)_{T}\left(\frac{\partial V}{\partial T}\right)_{p}}{\left(\frac{\partial V}{\partial p}\right)_{T}}
  \end{align}
\end{subequations}
which can be rearranged as
\begin{equation}
c_{p}-c_{V} = T \frac{\left(\frac{\partial S}{\partial p}\right)_{T}\left(\frac{\partial V}{\partial T}\right)_{p}}{\left(\frac{\partial V}{\partial p}\right)_{T}}.
\end{equation}
Then it is not too hard to manipulate the right-hand side to obtain
the desired form: the denominator is already equal to $V\kappa_{T}$,
the numerator is $T(\partial S/\partial p)_{T}V\alpha$ and then the
reciprocal of the Maxwell relation for internal energy should give the
desired result.
\end{proof}